\chapter{Fundamentos Pedag{\'o}gicos}
\label{cap:fundamentos}

\begin{objetivos}
\begin{itemize}
  \item Aplicar o modelo tridimensional de feedback (Feed Up, Feed Back, Feed Forward) na devolutiva de trabalhos acad{\^e}micos;
  \item Classificar objetivos de aprendizagem nos seis n{\'i}veis da Taxonomia Revisada de Bloom;
  \item Estruturar rubricas anal{\'i}ticas com r{\'o}tulos desenvolvimentais, incluindo o crit{\'e}rio de transpar{\^e}ncia no uso de IA;
  \item Relacionar o planejamento reverso (Backward Design) e o Desenho Universal para a Aprendizagem ao desenho de disciplinas.
\end{itemize}
\end{objetivos}

\section{O Feedback Formativo: Conceito e Mecanismos de A{\c c}{\~a}o}

O feedback {\'e} reconhecido pela literatura em educa{\c c}{\~a}o como uma das interven{\c c}{\~o}es pedag{\'o}gicas com maior pot{\^e}ncia para impulsionar o desempenho discente \cite{hattie2007power}. Todavia, a efici{\^e}ncia do feedback n{\~a}o {\'e} uniforme: h{\'a} modalidades que aceleram a aprendizagem, enquanto interven{\c c}{\~o}es mal estruturadas podem gerar desengajamento, ansiedade ou paralisia.

\subsection{O Modelo Tridimensional de Hattie \& Timperley}

Em seu estudo seminal, \citeonline{hattie2007power} estabeleceram que o feedback verdadeiramente formativo precisa responder, de modo cont{\'i}nuo e integrado, a tr{\^e}s quest{\~o}es centrais:

\begin{enumerate}
  \item \textbf{\emph{Feed Up} (Para onde vou?)}: Refere-se {\`a} clareza sobre as metas de aprendizagem, os crit{\'e}rios de sucesso e o padr{\~a}o de desempenho esperado. Se o estudante n{\~a}o sabe exatamente o que precisa demonstrar, qualquer devolu{\c c}{\~a}o torna-se incompreens{\'i}vel;
  \item \textbf{\emph{Feed Back} (Como estou indo?)}: Consiste na an{\'a}lise descritiva e objetiva da produ{\c c}{\~a}o atual do estudante em confronto direto com os crit{\'e}rios estabelecidos, evidenciando o que j{\'a} foi atingido e o que apresenta desconformidade;
  \item \textbf{\emph{Feed Forward} (O que fazer a seguir?)}: Constitui o componente mais frequentemente negligenciado nas salas de aula tradicionais. Representa o roteiro de a{\c c}{\~o}es imediatas, concretas e pratic{\'a}veis que o estudante deve executar para fechar a lacuna identificada e avan{\c c}ar em complexidade.
\end{enumerate}

\subsection{Os Quatro N{\'i}veis de Foco do Feedback}

Al{\'e}m das tr{\^e}s quest{\~o}es, \citeonline{hattie2007power} categorizam o feedback em quatro n{\'i}veis hier{\'a}rquicos de foco:

\begin{itemize}
  \item \textbf{Foco na Tarefa (FT)}: Informa se a resposta est{\'a} correta ou incorreta e aponta informa{\c c}{\~o}es factuais faltantes. {\'E} {\'u}til para iniciantes, mas limitado para transfer{\^e}ncia de conhecimento;
  \item \textbf{Foco no Processo (FP)}: Examina as estrat{\'e}gias cognitivas, os m{\'e}todos de resolu{\c c}{\~a}o, as rela{\c c}{\~o}es causais e os passos l{\'o}gicos adotados pelo estudante. Apresenta alta efic{\'a}cia para o desenvolvimento de racioc{\'i}nio profundo;
  \item \textbf{Foco na Autorregula{\c c}{\~a}o (FR)}: Estimula o monitoramento aut{\^o}nomo, a autocorre{\c c}{\~a}o, a gest{\~a}o do tempo e a confian{\c c}a reflexiva do aluno. Apresenta o mais duradouro impacto na forma{\c c}{\~a}o acad{\^e}mica e profissional;
  \item \textbf{Foco no \emph{Self} / Pessoa (FS)}: Consiste em elogios ou cr{\'i}ticas gen{\'e}ricas dirigidas {\`a} pessoa do estudante (ex.: ``Voc{\^e} {\'e} muito inteligente'' ou ``Voc{\^e} n{\~a}o tem voca{\c c}{\~a}o para c{\'a}lculo''). A literatura evidencia que o feedback no n{\'i}vel da pessoa possui impacto pedag{\'o}gico nulo ou negativo, pois desvia a aten{\c c}{\~a}o da tarefa e refor{\c c}a mentalidades fixas de intelig{\^e}ncia.
\end{itemize}

\begin{conceitochave}
\textbf{Princ{\'i}pio da Carga Cognitiva no Feedback}: Um feedback eficaz limita-se a um n{\'u}mero gerenci{\'a}vel de a{\c c}{\~o}es de melhoria (idealmente entre uma e duas a{\c c}{\~o}es priorit{\'a}rias). Relat{\'o}rios com dezenas de corre{\c c}{\~o}es simult{\^a}neas sobrecarregam a mem{\'o}ria de trabalho do aluno e resultam em abandono.
\end{conceitochave}

Complementarmente, \citeonline{nicol2006formative} salientam que o bom feedback deve estimular o di{\'a}logo e a autoavalia{\c c}{\~a}o cr{\'i}tica discente. Em nossa arquitetura, os agentes e as skills de feedback operam com base estrita nesses princ{\'i}pios: analisam processos, formulam perguntas ativadoras da metacogni{\c c}{\~a}o e estruturam passos de \emph{feed forward} claros e vi{\'a}veis.

\section{A Taxonomia Revisada de Bloom}

A classifica{\c c}{\~a}o de objetivos educacionais proposta originalmente por \citeonline{bloom1956taxonomy} foi substancialmente reformulada por \citeonline{anderson2001taxonomy}. A \textbf{Taxonomia Revisada de Bloom (TRB)} converteu os substantivos originais em formas verbais din{\^a}micas e estruturou o dom{\'i}nio cognitivo em uma matriz bidimensional: a \emph{Dimens{\~a}o do Processo Cognitivo} e a \emph{Dimens{\~a}o do Conhecimento}.

\subsection{A Dimens{\~a}o do Processo Cognitivo}

A progress{\~a}o cognitiva compreende seis n{\'i}veis em ordem crescente de complexidade, conforme sintetizado na Tabela~\ref{tab:bloom-niveis}.

\begin{table}[htbp]
  \centering
  \small
  \caption{Os 6 N{\'i}veis Cognitivos da Taxonomia Revisada de Bloom e o Impacto da IA}
  \label{tab:bloom-niveis}
  \begin{tabularx}{\linewidth}{clp{4.5cm}X}
    \toprule
    \textbf{N{\'i}vel} & \textbf{Processo} & \textbf{Verbos de A{\c c}{\~a}o} & \textbf{Vulnerabilidade / Oportunidade com IA} \\
    \midrule
    \textbf{1} & Lembrar & Listar, reconhecer, definir, nomear, reproduzir & \emph{Alta vulnerabilidade}. A IA responde instantaneamente; valor pedag{\'o}gico isolado reduzido. \\
    \addlinespace
    \textbf{2} & Entender & Explicar, exemplificar, resumir, interpretar, classificar & \emph{M{\'e}dia vulnerabilidade}. LLMs parafraseiam conceitos; exige exig{\^e}ncia de exemplos pr{\'o}prios e locais. \\
    \addlinespace
    \textbf{3} & Aplicar & Implementar, resolver, executar, calcular, usar & \emph{M{\'e}dia-alta}. A IA executa procedimentos padr{\~a}o; requer problemas contextualizados e dados novos. \\
    \addlinespace
    \textbf{4} & Analisar & Comparar, diferenciar, organizar, desconstruir, auditar & \emph{Alta oportunidade}. O aluno analisa erros em c{\'o}digos ou relat{\'o}rios gerados pela pr{\'o}pria IA. \\
    \addlinespace
    \textbf{5} & Avaliar & Julgar, criticar, defender, priorizar, validar & \emph{Elevada relev{\^a}ncia}. Demanda julgamento de valor, an{\'a}lise de vieses {\'e}ticos e conformidade normativa. \\
    \addlinespace
    \textbf{6} & Criar & Projetar, formular, planejar, sintetizar, construir & \emph{Parceria criativa}. Co-cria{\c c}{\~a}o de solu{\c c}{\~o}es originais onde o estudante assume a dire{\c c}{\~a}o autoral. \\
    \bottomrule
  \end{tabularx}
\end{table}

\subsection{A Dimens{\~a}o do Conhecimento}

Paralelamente aos processos mentais, os objetivos educacionais articulam quatro tipos de conhecimento \cite{anderson2001taxonomy}:
\begin{enumerate}
  \item \textbf{Conhecimento Factual}: Elementos b{\'a}sicos, terminologias e detalhes espec{\'i}ficos;
  \item \textbf{Conhecimento Conceitual}: Esquemas, categorias, princ{\'i}pios e generaliza{\c c}{\~o}es estruturantes;
  \item \textbf{Conhecimento Procedimental}: M{\'e}todos, t{\'e}cnicas, algoritmos e crit{\'e}rios para aplicar solu{\c c}{\~o}es;
  \item \textbf{Conhecimento Metacognitivo}: Consci{\^e}ncia da pr{\'o}pria cogni{\c c}{\~a}o, conhecimento estrat{\'e}gico e gest{\~a}o do pr{\'o}prio aprendizado.
\end{enumerate}

\begin{conceitochave}
\textbf{Independ{\^e}ncia entre Bloom e AIAS}: {\'E} fundamental n{\~a}o confundir o n{\'i}vel de Bloom com o n{\'i}vel AIAS. Uma atividade no n{\'i}vel AIAS~1 (Sem IA) pode demandar processos de ordem superior (\emph{Criar} ou \emph{Avaliar}), como uma prova pr{\'a}tica de bancada ou reda{\c c}{\~a}o reflexiva manuscrita. Inversamente, uma atividade AIAS~4 (IA Integral) pode focar no n{\'i}vel cognitivo \emph{Analisar}, exigindo que o aluno audite e encontre vulnerabilidades de seguran{\c c}a em um c{\'o}digo gerado por IA.
\end{conceitochave}

\section{Rubricas de Avalia{\c c}{\~a}o Formativa}

Uma rubrica {\'e} um guia pontuado e estruturado que combina crit{\'e}rios de avalia{\c c}{\~a}o com descritores qualitativos de desempenho para diferentes n{\'i}veis de dom{\'i}nio \cite{brookhart2013rubrics}. As rubricas cumprem papel decisivo na avalia{\c c}{\~a}o formativa:
\begin{itemize}
  \item Comunicam previamente as expectativas aos estudantes;
  \item Reduzem a subjetividade e o vi{\'e}s inconsciente na corre{\c c}{\~a}o;
  \item Possibilitam a autoavalia{\c c}{\~a}o e a avalia{\c c}{\~a}o por pares com rigor \cite{andrade2019critical};
  \item Transformam a corre{\c c}{\~a}o em um mapa claro de \emph{feed back} e \emph{feed forward}.
\end{itemize}

\subsection{Estrutura de uma Rubrica Anal{\'i}tica Robusta}

Conforme prescreve a skill \texttt{ia-educacao-rubrica}, as rubricas desenvolvidas no reposit{\'o}rio s{\~a}o prioritariamente \textbf{anal{\'i}ticas} (avaliam crit{\'e}rios individualmente) e seguem crit{\'e}rios rigorosos de desenho:

\begin{enumerate}
  \item \textbf{Espelhamento de Objetivos}: Cada crit{\'e}rio da rubrica deve espelhar diretamente um objetivo de aprendizagem declarado. Crit{\'e}rios cosm{\'e}ticos desprovidos de respaldo formativo (ex.: ``est{\'e}tica geral'') s{\~a}o desencorajados;
  \item \textbf{Extens{\~a}o Gerenci{\'a}vel}: Recomenda-se de 3 a 6 crit{\'e}rios por rubrica, com 4 n{\'i}veis de desempenho;
  \item \textbf{R{\'o}tulos Neutros e Desenvolvimentais}: Adotam-se r{\'o}tulos que valorizam o percurso de aprendizagem:
    \begin{itemize}
      \item \emph{Iniciante}: O estudante compreende os objetivos b{\'a}sicos, mas apresenta inconsist{\^e}ncias conceituais ou lacunas t{\'e}cnicas substanciais;
      \item \emph{Em desenvolvimento}: Demonstra compreens{\~a}o parcial e aplica procedimentos essenciais, necessitando de ajustes e refinamentos;
      \item \emph{Proficiente}: Atende plenamente aos objetivos de aprendizagem com rigor t{\'e}cnico, clareza e autonomia;
      \item \emph{Exemplar}: Supera as expectativas m{\'i}nimas, integrando an{\'a}lise cr{\'i}tica original, conex{\~o}es interdisciplinares e excel{\^e}ncia metodol{\'o}gica.
    \end{itemize}
  \item \textbf{Crit{\'e}rio de Coer{\^e}ncia com a Escala AIAS}: Em tarefas que admitem ou exigem ferramentas computacionais, a rubrica inclui um crit{\'e}rio espec{\'i}fico sobre transpar{\^e}ncia e integridade no uso de IA, avaliando se a autoria e o refinamento respeitaram o n{\'i}vel AIAS acordado.
\end{enumerate}

\section{Planejamento Did{\'a}tico e Inclus{\~a}o}

\subsection{Planejamento Reverso (\emph{Backward Design})}

A metodologia de \emph{Backward Design}, formulada por \citeonline{wiggins2005understanding}, estabelece que o planejamento did{\'a}tico deve come{\c c}ar n{\~a}o pela escolha de conte{\'u}dos ou listas de livros, mas pelo destino final da aprendizagem. A estrutura organiza-se em tr{\^e}s etapas sequenciais:

\begin{enumerate}
  \item \textbf{Etapa 1 -- Identificar os Resultados Desejados}: O que os estudantes devem ser capazes de saber, compreender e fazer ao t{\'e}rmino do percurso? Quais s{\~a}o as ideias perenes e compreens{\~o}es essenciais?
  \item \textbf{Etapa 2 -- Determinar as Evid{\^e}ncias Aceit{\'a}veis}: Como saberemos se os estudantes alcan{\c c}aram esses resultados? Quais tarefas aut{\^e}nticas, avalia{\c c}{\~o}es formativas e rubricas fornecer{\~a}o evid{\^e}ncias leg{\'i}timas de aprendizagem?
  \item \textbf{Etapa 3 -- Planejar as Experi{\^e}ncias de Aprendizagem e o Ensino}: Quais atividades ativas, leituras e interven{\c c}{\~o}es pedag{\'o}gicas conduzir{\~a}o os discentes {\`a} constru{\c c}{\~a}o das compet{\^e}ncias demandadas?
\end{enumerate}

Esse fluxo dialoga diretamente com o princ{\'i}pio do \textbf{Alinhamento Construtivo} de \citeonline{biggs2011teaching}, segundo o qual o processo de ensino apenas tem {\^e}xito quando as tarefas avaliativas e as estrat{\'e}gias did{\'a}ticas operam em perfeita sintonia cognitiva com os objetivos almejados.

\subsection{Desenho Universal para a Aprendizagem (DUA) e Acessibilidade}

Nenhum sistema educacional pode ser considerado de qualidade se n{\~a}o for inclusivo. O referencial do \textbf{Desenho Universal para a Aprendizagem (DUA)} \cite{cast2018udl} prop{\~o}e a elimina{\c c}{\~a}o sistem{\'a}tica de barreiras pedag{\'o}gicas por meio de tr{\^e}s eixos fundantes:

\begin{itemize}
  \item \textbf{M{\'u}ltiplos Meios de Engajamento} (o ``porqu{\^e}'' da aprendizagem): estimular o interesse, a persist{\^e}ncia e a autorregula{\c c}{\~a}o emocional;
  \item \textbf{M{\'u}ltiplos Meios de Representa{\c c}{\~a}o} (o ``o qu{\^e}'' da aprendizagem): apresentar conceitos por diferentes linguagens (textual, visual, gr{\'a}fica, auditiva e interativa);
  \item \textbf{M{\'u}ltiplos Meios de A{\c c}{\~a}o e Express{\~a}o} (o ``como'' da aprendizagem): oportunizar diferentes formas para que os estudantes demonstrem suas compet{\^e}ncias.
\end{itemize}

No {\^a}mbito de nosso acervo de agentes, a skill \texttt{ia-educacao-dua} e o agente \texttt{coach-de-feedback-formativo} aplicam princ{\'i}pios do DUA na produ{\c c}{\~a}o de feedbacks com linguagem concisa, organiza{\c c}{\~a}o em etapas e redu{\c c}{\~a}o de jarg{\~o}es desnecess{\'a}rios. As adapta{\c c}{\~o}es devem responder {\`a}s necessidades e prefer{\^e}ncias informadas pelo estudante, sem inferir diagn{\'o}sticos. A estrutura das skills (padr{\~a}o \texttt{SKILL.md}) e o cat{\'a}logo completo de subagentes e skills s{\~a}o apresentados nos Cap{\'i}tulos~\ref{cap:skills-opencode} e~\ref{cap:agentes-skills}.
